\documentclass[11pt]{article}
\usepackage[margin=1in]{geometry}
\usepackage{amsmath,amssymb}
\usepackage{array}
\usepackage{enumitem}
\usepackage{hyperref}
\usepackage[T1]{fontenc}

\setlength{\parindent}{0pt}
\setlength{\parskip}{0.6em}
\newcommand{\unchecked}{\(\square\)}

\title{Circuit Entanglement: Weekly Study Notes}
\author{Alexandra Bouillon}
\date{October 12--18, 2026}

\begin{document}
\maketitle

\section{Goal}

Understand how a quantum circuit creates a Bell pair, how to verify the result, and how that local circuit becomes the basic resource used by a quantum network.

\section{Core Circuit}

Start with two qubits in \( |00\rangle \):

\begin{verbatim}
q0: |0> ----H----o----
                 |
q1: |0> ---------X----
\end{verbatim}

\subsection{State after each step}

\begin{enumerate}
    \item Initial state:
    \[
    |q_0q_1\rangle=|00\rangle.
    \]

    \item Apply the Hadamard gate to \(q_0\):
    \[
    (H\otimes I)|00\rangle
    =\frac{|00\rangle+|10\rangle}{\sqrt{2}}.
    \]

    \item Apply CNOT with \(q_0\) as control and \(q_1\) as target:
    \[
    \frac{|00\rangle+|10\rangle}{\sqrt{2}}
    \longrightarrow
    \frac{|00\rangle+|11\rangle}{\sqrt{2}}
    =|\Phi^+\rangle.
    \]
\end{enumerate}

The qubits are now entangled: measuring one fixes the correlated outcome of the other. The individual result is random, but the pair is correlated.

\section{What Each Gate Contributes}

\begin{itemize}
    \item \textbf{Hadamard:} creates a superposition on the control qubit.
    \item \textbf{CNOT:} turns that superposition into a two-qubit correlation.
    \item \textbf{Measurement:} produces either \(00\) or \(11\) in the ideal computational-basis experiment.
\end{itemize}

The CNOT alone does not create this Bell state from \( |00\rangle \); it needs the superposition created by \(H\).

\section{Expected Ideal Measurements}

After adding computational-basis measurements:

\begin{verbatim}
q0: |0> ----H----o----M---- c0
                 |
q1: |0> ---------X----M---- c1
\end{verbatim}

Expected counts over many shots:

\begin{itemize}
    \item Approximately 50\% \(00\)
    \item Approximately 50\% \(11\)
    \item Ideally 0\% \(01\)
    \item Ideally 0\% \(10\)
\end{itemize}

Seeing only correlated computational-basis outcomes is necessary, but by itself it is not a complete proof of entanglement. A classical mixture can also produce \(00\) and \(11\). Verification in another basis or state tomography is needed to distinguish the Bell superposition from that classical mixture.

\section{Connection to a Physical Network Link}

The circuit above creates two entangled qubits locally. A network needs the halves at different nodes:

\begin{verbatim}
Bell-pair source
     |-- q0 --> Node A memory
     `-- q1 --> Node B memory
\end{verbatim}

Operational sequence:

\begin{enumerate}
    \item Create an entangled pair.
    \item Deliver or map one half to each adjacent node.
    \item Herald successful delivery through a classical message.
    \item Store both halves in quantum memory.
    \item Treat the pair as one available network ebit.
\end{enumerate}

The circuit describes the logical state preparation. Photon generation, loss, detection, memory mapping, and heralding are physical/link-layer operations around the circuit.

\section{Rate and Fidelity}

Keep these separate:

\begin{itemize}
    \item \textbf{Link success \(p\):} probability that a usable pair is established in the time window.
    \item \textbf{Fidelity \(F\):} quality of the pair after success.
\end{itemize}

For target \( |\Phi^+\rangle \) and actual state \(\rho\):

\[
F=\langle\Phi^+|\rho|\Phi^+\rangle.
\]

A link may have:

\begin{itemize}
    \item High rate and low fidelity: pairs arrive often but are noisy.
    \item Low rate and high fidelity: good pairs are rare.
    \item Low rate and low fidelity: both delivery and quality are poor.
\end{itemize}

\section{Working Checklist}

\begin{itemize}[leftmargin=*]
    \item[\unchecked] Draw the Bell circuit from memory.
    \item[\unchecked] Explain the state after \(H\).
    \item[\unchecked] Explain the state after CNOT.
    \item[\unchecked] Predict ideal measurement counts.
    \item[\unchecked] Explain why \(00/11\) counts alone do not fully prove entanglement.
    \item[\unchecked] Distinguish a local Bell circuit from remote entanglement distribution.
    \item[\unchecked] Explain heralding in one sentence.
    \item[\unchecked] Explain rate versus fidelity in one sentence each.
    \item[\unchecked] Build the circuit in Qiskit or another simulator.
    \item[\unchecked] Record ideal counts.
    \item[\unchecked] Add a simple noise model and record the change.
\end{itemize}

\section{Connection to the Current Labs}

\begin{itemize}
    \item \textbf{Labs 1--8:} polarization gives single-qubit state and gate intuition.
    \item \textbf{Lab 9:} two Bell links become the input resource for entanglement swapping.
    \item \textbf{Lab 10:} the Bell state cannot be copied and broadcast to arbitrary users.
    \item \textbf{Lab 11:} each Bell-link attempt has a success probability used in the route product.
\end{itemize}

\section{Questions to Answer This Week}

\begin{enumerate}
    \item Why does \(H\) have to come before CNOT?
    \item What changes if the control and target are reversed?
    \item How would I prepare \( |\Phi^-\rangle \), \( |\Psi^+\rangle \), and \( |\Psi^-\rangle \)?
    \item What experiment distinguishes \( |\Phi^+\rangle \) from a classical 50/50 mixture of \( |00\rangle \) and \( |11\rangle \)?
    \item Where does the pair source belong when Node A and Node B are physically separated?
\end{enumerate}

\section{My Explanation}

Write a three-sentence explanation without looking above:

\begin{enumerate}
    \item \vspace{1cm}
    \item \vspace{1cm}
    \item \vspace{1cm}
\end{enumerate}

\end{document}
